
\section*{Conclusions}
\addcontentsline{toc}{section}{Conclusions}

Mixture of Experts architectures represent a paradigm shift in how we design and deploy AI systems for healthcare. By enabling conditional computation through specialized expert pathways and learned routing mechanisms, MoE addresses fundamental challenges that have limited the effectiveness of conventional deep learning in clinical settings: data heterogeneity, multi-modality, task diversity, and the need for scalable yet efficient models.

The applications surveyed in this Review span the full breadth of healthcare AI, from medical imaging and computational pathology to clinical NLP, drug discovery, EHR analysis, and multi-modal clinical decision support. Across these domains, a consistent theme emerges: the modular, conditional nature of MoE architectures aligns naturally with the structure of medical data and clinical reasoning. Diseases manifest differently across patients and modalities, clinical expertise is inherently specialized, and effective healthcare delivery requires the integration of diverse knowledge sources. MoE provides a computational framework that mirrors these characteristics, enabling AI systems that are more adaptable, more efficient, and more aligned with clinical practice than their monolithic counterparts.

However, significant challenges remain before MoE-based clinical AI can achieve widespread deployment. Interpretability requirements demand that expert routing and expert computations be transparent and clinically meaningful. Expert collapse must be prevented to maintain the benefits of specialization, particularly for rare conditions and underrepresented populations. Data privacy constraints require the development of federated MoE training approaches that preserve patient confidentiality while enabling multi-institutional collaboration. Regulatory frameworks must evolve to address the unique characteristics of modular, conditionally activated AI systems.

Looking forward, the convergence of MoE with emerging trends in clinical AI, including foundation models, precision medicine, and causal inference, promises transformative advances. Personalized expert routing, dynamic expert creation, and the integration of domain knowledge into gating mechanisms represent particularly compelling directions. As these technologies mature, MoE architectures have the potential to become a foundational building block of clinical AI systems that are not only more powerful but also more interpretable, more equitable, and more aligned with the collaborative, specialized nature of clinical medicine.

The journey from the original MoE formulation three decades ago to its current applications in healthcare illustrates a broader lesson in AI research: architectures that respect the structure of the problem domain consistently outperform those that do not. In healthcare, where the stakes are highest and the data most complex, this principle is more important than ever. Mixture of Experts, by embracing the inherent modularity and specialization of medical knowledge, offers a path toward AI systems that truly augment rather than replace clinical expertise.
